\scanpage{202}
\mspage{-2.12-}
% archive sheet number: 201
The ``only if'' in (ii) is proven in the same way: assume $(0)\handwritten{\unsure{\in|\epsilon}}$ closure $O(x^*)$. As a first step, we obtain a $K$-valued point $\handwritten{\phi}$ of $G$ such that $(0)$ is the specialization of $\handwritten{\psi_{\unsure{x|\chi}^*}\circ\phi}$ when $t\handwritten{\longrightarrow}0$. We decompose $\handwritten{\phi}$ exactly as before, and obtain $(*)\handwritten{'}$. But now the term on the left hand side of $(*)\handwritten{'}$ is actually in $t\cdot R$ since $\handwritten{\sigma}^{*}\handwritten{\bigl(}(\handwritten{\psi}_1 b_2)^{-1},\handwritten{\psi_{x^*}\circ\phi\bigr)}$ actually specializes to $(0)$ when $t\handwritten{\longrightarrow}0$. Therefore, instead of (A), we obtain:
\[
X_i\handwritten{\bigl\{}\handwritten{\sigma}^{*}(b_2^{-1}\handwritten{\psi}_2,\ x^*)\handwritten{\bigr\}}\handwritten{\in}t^{-r_{\handwritten{i}}+1}\cdot R\ .
\tag*{($\overset{\handwritten{\sim}}{A}$)}
\]

(B) is the same as before, hence combining $(\overset{\handwritten{\sim}}{A})$ and (B) it follows that $r_i>0$ whenever $X_i(x^*)\handwritten{\neq}0$. Therefore, $\handwritten{\mu}(x,b_2^{-1}\handwritten{\lambda}\ b_2)<0$ and (ii) is proven. \uline{QED}

A second application of Iwahori's theorem is:

\uline{Proposition} \uline{2.4}: Let a reductive algebraic group $G$ over $k$ act via $\handwritten{\sigma}$ on an algebraic scheme $X$ over $k$. Then the action of $G$ is proper if and only if for every non-trivial 1-PS $\handwritten{\lambda}:\handwritten{\mathbb{G}}_m\handwritten{\longrightarrow}G$, the induced action of $\handwritten{\mathbb{G}}_m$ on $X$ is proper.
% vertical marginal rule beside Proposition 2.4
\figure{\begin{tikzpicture}\draw (0,0)--(0,1.8);\end{tikzpicture}}

\uline{Proof}: The only if is clear. To prove the if, suppose the action of $G$ is not proper, i.e. $\handwritten{\Psi}:G\times X\handwritten{\longrightarrow}X\times X$ is not proper. Then, by the valuative criterion for properness, there are $K$-valued points $\handwritten{\phi}$ and $\handwritten{\xi}$ of $G$ and $X$ respectively such that $\handwritten{\Psi(\phi\times\xi)}$ is an $R$-valued point of $X\times X$, but $\handwritten{\phi\times\xi}$ is not an $R$-valued point of $G\times X$. In other words, $\handwritten{\xi}$ and $\handwritten{\sigma(\phi,\xi)}$ are $R$-valued points of $X$, but $\handwritten{\phi}$ is not an $R$-valued point of $G$. By Iwahori's theorem, we can write
